%% kapa: exam=1 round=24 year=2013 period=1 form=A subject=economics subject_ko=경제원론 seq=39 number=79
%% source: Q-Net 감정평가사 기출문제 (cst00309) artlSeq=5073802 file=extracted/1008/5073802/2013년 제24회 감정평가사 1차시험 기출문제/1교시(A형).hwp
%% conversion: kordoc 4.24.51 -> tools/gen1.py (auto); review=proofread-sonnet; flags=none
\begin{question}{79}
중앙은행은 확장적 통화정책을 실시하여 균형국민소득을 60만큼 더 증가시키고자 한다. 경제모형이 다음과 같을 때 중앙은행은 통화량을 현재의 수준에서 얼마만큼 더 증가시켜야 하는가? (단, $Y$ 는 국민소득, $C$ 는 소비, $I$ 는 투자, $G$ 는 정부지출, $r$ 은 이자율, $M^{ d }$ 는 화폐수요, $M^{ s }$ 는 화폐공급이며, 물가수준은 고정되어 있다.)
\begin{framedtext}
$C=10+0.7Y$ \par
$I=50-r$ \par
$G=10$ \par
$Y=C+I+G$ (생산물시장의 균형) \par
$M^{ d } =20+0.2Y-r$ \par
$M^{ s } =30$ \par
$M^{ d } =M^{ s }$ (화폐시장의 균형)
\end{framedtext}
\begin{choices}
  \item 10
  \item 20
  \item 30
  \item 40
  \item 50
\end{choices}
\end{question}
\answer{3}
